\documentclass[11pt]{article}
\usepackage[a4paper,margin=1in]{geometry}
\usepackage{amsmath,amssymb,mathtools}
\usepackage{enumitem}

\newcommand{\IC}{\operatorname{IC}}
\newcommand{\PR}{\operatorname{PR}}
\newcommand{\rank}{\operatorname{rank}}
\newcommand{\sgn}{\operatorname{sgn}}
\newcommand{\finite}{\operatorname{finite}}
\newcommand{\corr}{\operatorname{corr}}
\newcommand{\Var}{\operatorname{Var}}
\newcommand{\E}{\operatorname{E}}
\newcommand{\ind}{\mathbf{1}}

\title{Factor Research Evolution: Mathematical Summary}
\date{}

\begin{document}
\maketitle

\section*{Core Objects}

Let \(P_{t,a}\) be close price for time \(t=1,\dots,T\) and asset
\(a\in\mathcal A\).  A candidate factor signal is \(x_{t,a}\); the current
accepted book is \(B=\{b_1,\dots,b_m\}\).  The target horizon is
\(h=\texttt{cfg.target\_horizon}\).  Forward returns are
\[
  r^{(h)}_{t,a}=\frac{P_{t+h,a}}{P_{t,a}}-1.
\]
The temporal split is \(I\) = in-sample, \(V\) = validation, \(H\) = final test,
and \(D=I\cup V\).  The search fitness uses \(I\) and \(V\); \(H\) is not used
inside \(\texttt{evaluate\_candidate}\) or \(\texttt{evaluate\_set}\).

For any row mask \(M\), the label-availability mask used by the code is
\[
  L_h(M)_t = M_t \land M_{t+h},
\]
with out-of-range indices treated as false.  This prevents validation labels
from consuming final-test prices.

Per-underlying z-scoring over statistic rows \(R\) is
\[
  z^R(x)_{t,a}=\frac{x_{t,a}-\mu^R_a(x)}{\sigma^R_a(x)},\qquad
  \mu^R_a(x)=\frac{1}{|R_a|}\sum_{u\in R}x_{u,a},
\]
with \(\sigma^R_a(x)\) the sample standard deviation; zero standard deviations
and infinities become missing.

The pooled Pearson IC used throughout is
\[
  \IC_h(x;M,R,A)
  =
  \rho\!\left(
    \operatorname{vec}_{t\in M\cap L_h(A),a}\, z^R(x)_{t,a},
    \operatorname{vec}_{t\in M\cap L_h(A),a}\, r^{(h)}_{t,a}
  \right),
\]
where for finite paired vectors \(u,v\),
\[
  \rho(u,v)=
  \frac{\sum_i(u_i-\overline u)(v_i-\overline v)}
       {\sqrt{\sum_i(u_i-\overline u)^2}
        \sqrt{\sum_i(v_i-\overline v)^2}},
\]
and the result is undefined if fewer than three finite pairs exist.

For a book \(B\), the combined predictor \(\widehat y_B\) is obtained by fitting
the configured estimator \(\mathcal M\) on IS rows:
\[
  \widehat f_B
  =
  \mathcal M\!\left(
    \{z^I(b_j)_{t,a}\}_{j=1}^m,\ r^{(h)}_{t,a}
  \right)_{t\in I\cap L_h(I),a},
  \qquad
  \widehat y_B(t,a)=\widehat f_B\!\left(z^I(b_1)_{t,a},\dots,z^I(b_m)_{t,a}\right).
\]
The default \(\mathcal M\) is \(\texttt{gradient\_boosting}\).  Missing and
infinite features are imputed to zero before fitting/prediction.

\section*{SINGLE-Mode Reward Vector}

For one candidate \(x\), define
\[
  \IC_{\mathrm{IS}}(x)=\IC_h(x;I,I,I),\qquad
  \IC_{\mathrm{VAL}}(x)=\IC_h(x;V,I,D).
\]
The Pareto objective is
\[
  F(x)=
  \bigl(F_1,F_2,F_3,F_4,F_5\bigr),
  \qquad \text{all axes are maximized}.
\]

\paragraph{1. Marginal value.}
Let \(B_x=B\cup\{x\}\).  If \(B\neq\varnothing\),
\[
  F_1(x)
  =
  \IC_h(\widehat y_{B_x};V,I,D)
  -
  \IC_h(\widehat y_B;V,I,D).
\]
If \(B=\varnothing\), the base IC is set to \(0\), so
\[
  F_1(x)=\IC_h(\widehat y_{\{x\}};V,I,D).
\]

\paragraph{2. Independence, default: residual IC.}
Let \(c=\operatorname{vec} z^I(x)\) and
\[
  Z_B=\begin{bmatrix}
  \operatorname{vec} z^I(b_1) & \cdots & \operatorname{vec} z^I(b_m)
  \end{bmatrix}.
\]
For \(m>0\), fit on finite IS cells
\[
  c=\alpha\mathbf 1 + Z_B\beta+\varepsilon.
\]
Then
\[
  F_2(x)=
  \rho\!\left(
    \varepsilon_{t\in V\cap L_h(D),a},
    r^{(h)}_{t\in V\cap L_h(D),a}
  \right).
\]
If \(B=\varnothing\), this reduces to
\[
  F_2(x)=\IC_h(x;V,I,D).
\]

\paragraph{2b. Optional legacy independence: delta participation.}
When \(\texttt{independence\_metric}=\texttt{delta\_participation}\),
\[
  \PR(C)=\frac{\left(\sum_i\lambda_i(C)\right)^2}{\sum_i\lambda_i(C)^2},
\]
where \(\lambda_i(C)\) are clipped nonnegative eigenvalues of the signal
correlation matrix \(C\).  Using flattened \(z^D(\cdot)\) over \(D\),
\[
  \Delta\PR=\PR(C_{B\cup x})-\PR(C_B),\qquad
  \rho_{\max}=\max_{b\in B}\left|\corr(z^D(x),z^D(b))\right|,
\]
\[
  F_2(x)=\Delta\PR-\lambda_{\mathrm{corr}}\rho_{\max},
  \qquad \lambda_{\mathrm{corr}}=0.5.
\]
For an empty book the implementation returns \(F_2=1\).

\paragraph{3. Robustness.}
Let
\[
  m_x=F_1(x),\qquad
  s_m=\begin{cases}
    +1, & m_x\text{ is undefined or }m_x\ge 0,\\
    -1, & m_x<0.
  \end{cases}
\]
CPCV partitions \(D\) into \(G\) contiguous groups and tests every
\(k\)-group combination; defaults are \(G=6,k=2\), so \(\binom{6}{2}=15\)
folds.  For fold \(f\) with purged/embargoed train mask \(A_f\) and test mask
\(T_f\), fit fresh combined predictors on \(A_f\):
\[
  \widehat y^{(f)}_{B_x}
  =
  \mathcal M(B\cup\{x\};A_f),
  \qquad
  \widehat y^{(f)}_{B}
  =
  \mathcal M(B;A_f).
\]
The scored CPCV value is fold-refit LOCO marginal IC:
\[
  q_f
  =
  s_m\left[
    \IC_h(\widehat y^{(f)}_{B_x};T_f,A_f,D)
    -
    \IC_h(\widehat y^{(f)}_{B};T_f,A_f,D)
  \right],
\]
with base \(0\) when \(B=\varnothing\).  The previous standalone CPCV stability
probe is still reported as a diagnostic:
\[
  u_f=s_x\,\IC_h(x;T_f,A_f,D),\qquad
  s_x=\sgn_{\mathrm{impl}}(\IC_{\mathrm{VAL}}(x)).
\]
If at least two \(q_f\) values exist,
\[
  \bar q=\frac{1}{|\mathcal F|}\sum_{f\in\mathcal F}q_f,\qquad
  \sigma_q=\sqrt{\frac{1}{|\mathcal F|}\sum_{f\in\mathcal F}(q_f-\bar q)^2}.
\]
For window-jitter variants \(x^{(j)}\),
\[
  P_{\mathrm{plateau}}
  =
  \max\left(
    0,\,
    s_x\IC_{\mathrm{VAL}}(x)
    -
    \frac{1}{J}\sum_{j=1}^J s_x\IC_h(x^{(j)};V,I,D)
  \right).
\]
If perturbation is enabled, with deterministic Gaussian shocks
\(\xi^{(d)}_{t,a}\sim N(0,\sigma_{\xi}^2\,\hat\sigma^2_{D,a}(x))\),
\[
  P_{\xi}
  =
  \max\left(
    0,\,
    s_x\IC_{\mathrm{VAL}}(x)
    -
    \frac{1}{D_\xi}\sum_{d=1}^{D_\xi}
      s_x\IC_h(x+\xi^{(d)};V,I,D)
  \right).
\]
The sign bonus is
\[
  B_{\mathrm{sign}}
  =
  \begin{cases}
    +b, & \sgn_{\mathrm{impl}}(\IC_{\mathrm{VAL}}(x))=\sgn(e),\\
    -b, & \sgn_{\mathrm{impl}}(\IC_{\mathrm{VAL}}(x))\neq\sgn(e),\\
    0, & e\text{ or }\IC_{\mathrm{VAL}}\text{ is undefined},
  \end{cases}
  \qquad b=0.02,
\]
where \(e\) is the expected sign and
\(\sgn_{\mathrm{impl}}(u)=+1\) for \(u\ge 0\), else \(-1\).  Thus
\[
  F_3(x)
  =
  \bar q-\lambda_{\mathrm{std}}\sigma_q
  -w_pP_{\mathrm{plateau}}
  -w_\xi P_\xi
  +B_{\mathrm{sign}}.
\]
Defaults: \(\lambda_{\mathrm{std}}=1\), \(w_p=1\), \(w_\xi=0\).

\paragraph{4. Parsimony.}
Let \(C(\text{code})\) count AST complexity:
binary ops, unary ops, calls, numeric constants each count \(1\);
a boolean op counts \(\max(1,n_{\mathrm{values}}-1)\);
a comparison counts its number of comparison operators;
syntax error gives \(10{,}000\).  Then
\[
  F_4(x)=-C(\text{code}_x).
\]

\paragraph{5. Regime independence.}
For default \(\texttt{regime\_kind}=\texttt{drawdown}\), define market return
\[
  m_t=\frac{1}{|\mathcal A|}\sum_{a\in\mathcal A}\left(\frac{P_{t,a}}{P_{t-1,a}}-1\right),
\]
and stress rows
\[
  S_t=D_t\land
  m_t\le Q_{\alpha}\left(\{m_u:u\in D,\ m_u\ \mathrm{finite}\}\right),
  \qquad \alpha=0.2.
\]
For \(\texttt{regime\_kind}=\texttt{volatility}\), \(m_t\) is replaced by a rolling
standard deviation \(v_t=\operatorname{sd}(m_{t-w+1:t})\), \(w=20\), and
\(S_t=D_t\land v_t\ge Q_{1-\alpha}(v)\).  The regime axis is
\[
  F_5(x)
  =
  \IC_h(\widehat y_{B_x};V\cap S,I,D)
  -
  \IC_h(\widehat y_B;V\cap S,I,D),
\]
with base \(0\) when \(B=\varnothing\).  The axis is undefined unless at least
\(\texttt{regime\_min\_obs}=20\) scored cells exist.

\section*{SET-Mode Reward Vector}

For a genome containing member signals \(M=\{x_1,\dots,x_K\}\), fit the combined
predictor \(\widehat y_M\) on those members only.  The five axes are
\[
  F_1(M)=\IC_h(\widehat y_M;V,I,D),
\]
\[
  F_2(M)=
  \begin{cases}
    \PR(C_M)/K, & K\ge 2,\\
    1, & K=1,
  \end{cases}
\]
where \(C_M\) is the correlation matrix of flattened \(z^D(x_i)\);
\[
  F_3(M)=\text{the same fold-refit robustness formula applied to }M
  \text{ as a self-contained combined model},
\]
\[
  F_4(M)=
  -\frac{1}{K}\sum_{i=1}^{K} C(\text{code}_{x_i}),
\]
\[
  F_5(M)=\IC_h(\widehat y_M;V\cap S,I,D).
\]
SET mode does not compute LOCO against an external archive; the set is scored as
a self-contained alpha program.

\section*{Hard Gates and Diagnostics}

Coverage over \(D\) is
\[
  \operatorname{cov}(x)=
  \frac{\sum_{t\in D,a}\ind[\finite(x_{t,a})]}
       {\sum_{t\in D,a}1}.
\]
The coverage gate is
\[
  g_{\mathrm{cov}}=\ind[\operatorname{cov}(x)\ge \tau_{\mathrm{cov}}],
  \qquad \tau_{\mathrm{cov}}=0.5.
\]

The degradation ratio is evaluated only if
\(\IC_{\mathrm{IS}},\IC_{\mathrm{VAL}}\) exist and
\(|\IC_{\mathrm{IS}}|\ge\texttt{min\_is\_ic}=0.005\):
\[
  d(x)=
  \frac{\IC_{\mathrm{VAL}}(x)\,\sgn(\IC_{\mathrm{IS}}(x))}
       {|\IC_{\mathrm{IS}}(x)|},
  \qquad
  g_{\mathrm{deg}}=\ind[d(x)\ge \tau_{\mathrm{deg}}],
  \qquad \tau_{\mathrm{deg}}=0.5.
\]
If the ratio is not evaluable, \(g_{\mathrm{deg}}\) is \(\texttt{None}\), not a
failure.

Deflation is currently diagnostic during search, not a per-candidate hard gate:
\[
  \texttt{deflation\_ok}=\texttt{None}.
\]
For \(n\) scored validation observations and \(N=\texttt{n\_trials}\),
\[
  \ell_{\IC}(n,N)=
  \begin{cases}
    0, & N\le 1 \text{ or } n\le 1,\\
    \sqrt{\frac{2\ln N}{n}}, & \text{otherwise},
  \end{cases}
\]
\[
  \IC_{\mathrm{deflated}}
  =
  \max\left(|\IC_{\mathrm{VAL}}|-\ell_{\IC}(n,N),0\right),
  \qquad
  t_{\mathrm{deflated}}
  =
  |\IC_{\mathrm{VAL}}|\sqrt n-\sqrt{2\ln N}.
\]
At publish time, when \(\texttt{selection\_deflation}=\texttt{on}\), the same
formula is applied to the final book's combined validation IC, and the book
passes iff \(t_{\mathrm{deflated}}>0\).

The optional turnover gate uses the explicit position primitive.  With default
threshold mode and causal expanding z-score,
\[
  p_{t,a}=
  \begin{cases}
    +1, & \tilde z_{t,a}>\theta,\\
    -1, & \tilde z_{t,a}<-\theta,\\
    0, & \text{otherwise},
  \end{cases}
  \qquad \theta=1.
\]
Then
\[
  \operatorname{turnover}(x)=
  \frac{1}{|D||\mathcal A|}
  \sum_{t\in D,a}|p_{t,a}-p_{t-1,a}|,
\]
\[
  \operatorname{gross}
  =
  \operatorname{mean}_{t\in D\cap L_h(D),a}(p_{t,a}r^{(h)}_{t,a}),
  \qquad
  \operatorname{net}
  =
  \operatorname{mean}_{t\in D\cap L_h(D),a}
  \left(p_{t,a}r^{(h)}_{t,a}-c|p_{t,a}-p_{t-1,a}|\right),
\]
with \(c=\texttt{cost\_rate}=5\cdot 10^{-4}\).  If
\(\tau_{\mathrm{turn}}\) is configured,
\[
  g_{\mathrm{cost}}=\ind[\operatorname{turnover}(x)\le \tau_{\mathrm{turn}}],
\]
otherwise \(g_{\mathrm{cost}}=\texttt{None}\).

A candidate is search-selectable iff every evaluated gate is true:
\[
  \operatorname{selectable}(x)
  =
  \bigwedge_{g\in\{g_{\mathrm{cov}},g_{\mathrm{deg}},g_{\mathrm{cost}}\},\ g\neq\texttt{None}} g.
\]

\section*{Selection: NSGA-II and QD}

For objective vectors \(F(x),F(y)\), treating undefined axes as \(-\infty\),
\[
  x\succ_P y
  \iff
  \left[\forall j,\ F_j(x)\ge F_j(y)\right]
  \land
  \left[\exists j,\ F_j(x)>F_j(y)\right].
\]
Constrained dominance is
\[
  x\succ_C y
  \iff
  \begin{cases}
    \text{true}, & \operatorname{selectable}(x)\land\neg\operatorname{selectable}(y),\\
    \text{false}, & \neg\operatorname{selectable}(x)\land\operatorname{selectable}(y),\\
    n_{\mathrm{fail}}(x)<n_{\mathrm{fail}}(y),
      & \neg\operatorname{selectable}(x),\neg\operatorname{selectable}(y),
        n_{\mathrm{fail}}(x)\ne n_{\mathrm{fail}}(y),\\
    x\succ_P y, & \text{otherwise}.
  \end{cases}
\]
The NSGA population is sorted by constrained nondominated fronts.  Within one
front, crowding distance for candidate \(i\) is
\[
  CD_i=\sum_{j=1}^{5}
  \frac{F_j(i_j^+)-F_j(i_j^-)}
       {\max_{\ell\in\mathrm{front}}F_j(\ell)-\min_{\ell\in\mathrm{front}}F_j(\ell)},
\]
with boundary points assigned \(CD_i=\infty\).  Parent selection is binary
tournament on lower front rank, then higher crowding distance.

In QD mode the objective vector is unchanged.  The behavior descriptors are
\[
  d_1(x)=
  \rho\!\left(
    \operatorname{vec}_{t\in D,a}z^D(x)_{t,a},
    \operatorname{vec}_{t\in D,a}\left(\frac{P_{t,a}}{P_{t-h,a}}-1\right)
  \right),
\]
\[
  d_2(x)=
  1-
  \frac{1}{|\mathcal A'|}
  \sum_{a\in\mathcal A'}
  \corr\!\left(z^D(x)_{1:T-1,a},z^D(x)_{2:T,a}\right),
\]
\[
  d_3(x)=
  \frac{\operatorname{mean}_{t\in D\cap S,a}|z^D(x)_{t,a}|}
       {\operatorname{mean}_{t\in D\setminus S,a}|z^D(x)_{t,a}|}.
\]
They are bucketed by
\[
  \operatorname{bucket}(v;e_1,\dots,e_q)=\sum_{r=1}^{q}\ind[v\ge e_r],
\]
with edges
\[
  d_1:\ (-0.33,0.33),\qquad
  d_2:\ (0.5,1.0),\qquad
  d_3:\ (0.9,1.3).
\]
The default QD cell is
\[
  \left(\operatorname{bucket}(d_1),\operatorname{bucket}(d_2)\right),
\]
and 3D QD adds \(\operatorname{bucket}(d_3)\).  Only selectable candidates enter
cells.  Each cell stores at most \(K=\texttt{cell\_capacity}=3\) elites, pruned by
the same constrained NSGA ranking.  Parent sampling chooses an occupied cell
uniformly, then an elite uniformly unless the optional weight
\[
  w_i=(1-\gamma)^{\operatorname{depth}_i}(1-\omega)^{\operatorname{reuse}_i}
\]
is enabled.  The reported QD score proxy is
\[
  \operatorname{QDScore}=\sum_{\mathrm{cell}}\max_{x\in\mathrm{cell}} F_1(x).
\]

\section*{End-of-Run Book Curation and Publish Filter}

Default persistence uses \(\texttt{accepted\_book}\): the Pareto archive for NSGA,
or the union of QD cell elites for QD.  If curation is enabled, it instead starts
from \(\texttt{kept\_pool}\), every gate-passing genome ever evaluated.

Greedy curation forward-selects the factor with largest combined validation IC:
\[
  x^*=\arg\max_{x\in R}\IC_h(\widehat y_{S\cup\{x\}};V,I,D),
  \qquad
  \Delta=\IC_h(\widehat y_{S\cup\{x^*\}};V,I,D)-\IC_h(\widehat y_S;V,I,D).
\]
It stops when the requested size is reached or, if size is free, once
\(\Delta\le\texttt{min\_gain}\) after at least one factor is selected.

Elastic-net curation fits repeated subsample elastic-net models on development
rows \(D\cap L_h(D)\).  With \(R\) successful resamples,
\[
  \pi_j=\frac{1}{R}\sum_{r=1}^{R}\ind[|\widehat\beta^{(r)}_j|>10^{-8}],
\]
and keeps either the top \(\texttt{n\_keep}\) factors by \(\pi_j\), or all
factors with \(\pi_j\ge 0.6\), falling back to the top one.

Publish-time deflation, when on, repeatedly drops the member with smallest LOCO
contribution
\[
  \Delta_j=
  \IC_h(\widehat y_K;V,I,D)
  -
  \IC_h(\widehat y_{K\setminus\{j\}};V,I,D)
\]
until the final combined book has \(t_{\mathrm{deflated}}>0\), reaches the
minimum size, or every remaining member has positive marginal contribution.

\section*{Search Budget}

The evolution loop is generation-capped, not quality-targeted:
\[
  N_{\mathrm{attempts}}
  \le
  N_{\mathrm{seed\ evals}}
  +
  \texttt{generations}\cdot\texttt{children\_per\_generation}.
\]
\(\texttt{n\_trials}\) increments only for successful, non-duplicate evaluations.
Default CLI values are
\[
  \texttt{generations}=5,\quad
  \texttt{population}=10,\quad
  \texttt{children\_per\_gen}=8,\quad
  \texttt{seed\_ideas}=8,\quad
  \texttt{islands}=1.
\]

\end{document}
